\documentclass[10pt,a4paper]{amsart}
\usepackage{amsmath,amssymb,amsthm,graphicx}
\usepackage{verbatim}
\newenvironment{reference}{\comment}{\endcomment}
%\newenvironment{reference}{}{}
\newenvironment{slogan}{\comment}{\endcomment}
\newenvironment{history}{\comment}{\endcomment}

% For commutative diagrams we use Xy-pic
\usepackage[all]{xy}

% We use 2cell for 2-commutative diagrams.
\xyoption{2cell}
% Theorem environments.
%
\theoremstyle{plain}
\newtheorem{theorem}[subsection]{Theorem}
\newtheorem{proposition}[subsection]{Proposition}
\newtheorem{lemma}[subsection]{Lemma}

\theoremstyle{definition}
\newtheorem{definition}[subsection]{Definition}
\newtheorem{example}[subsection]{Example}
\newtheorem{exercise}[subsection]{Exercise}
\newtheorem{situation}[subsection]{Situation}

\theoremstyle{remark}
\newtheorem{remark}[subsection]{Remark}
\newtheorem{remarks}[subsection]{Remarks}

\numberwithin{equation}{subsection}

% Macros
%
\def\lim{\mathop{\mathrm{lim}}\nolimits}
\def\colim{\mathop{\mathrm{colim}}\nolimits}
\def\Spec{\mathop{\mathrm{Spec}}}
\def\Hom{\mathop{\mathrm{Hom}}\nolimits}
\def\Ext{\mathop{\mathrm{Ext}}\nolimits}
\def\SheafHom{\mathop{\mathcal{H}\!\mathit{om}}\nolimits}
\def\SheafExt{\mathop{\mathcal{E}\!\mathit{xt}}\nolimits}
\def\Sch{\mathit{Sch}}
\def\Mor{\mathop{\mathrm{Mor}}\nolimits}
\def\Ob{\mathop{\mathrm{Ob}}\nolimits}
\def\Sh{\mathop{\mathit{Sh}}\nolimits}
\def\NL{\mathop{N\!L}\nolimits}
\def\CH{\mathop{\mathrm{CH}}\nolimits}
\def\proetale{{pro\text{-}\acute{e}tale}}
\def\etale{{\acute{e}tale}}
\def\QCoh{\mathit{QCoh}}
\def\Ker{\mathop{\mathrm{Ker}}}
\def\Im{\mathop{\mathrm{Im}}}
\def\Coker{\mathop{\mathrm{Coker}}}
\def\Coim{\mathop{\mathrm{Coim}}}

% Boxtimes
%
\DeclareMathSymbol{\boxtimes}{\mathbin}{AMSa}{"02}

%
% Macros for moduli stacks/spaces
%
\def\QCohstack{\mathcal{QC}\!\mathit{oh}}
\def\Cohstack{\mathcal{C}\!\mathit{oh}}
\def\Spacesstack{\mathcal{S}\!\mathit{paces}}
\def\Quotfunctor{\mathrm{Quot}}
\def\Hilbfunctor{\mathrm{Hilb}}
\def\Curvesstack{\mathcal{C}\!\mathit{urves}}
\def\Polarizedstack{\mathcal{P}\!\mathit{olarized}}
\def\Complexesstack{\mathcal{C}\!\mathit{omplexes}}
% \Pic is the operator that assigns to X its picard group, usage \Pic(X)
% \Picardstack_{X/B} denotes the Picard stack of X over B
% \Picardfunctor_{X/B} denotes the Picard functor of X over B
\def\Pic{\mathop{\mathrm{Pic}}\nolimits}
\def\Picardstack{\mathcal{P}\!\mathit{ic}}
\def\Picardfunctor{\mathrm{Pic}}
\def\Deformationcategory{\mathcal{D}\!\mathit{ef}}
\usepackage{cleveref}
\crefname{originallemma}{original lemma}{original lemmas}
\Crefname{originallemma}{Original lemma}{Original lemmas}
\makeatletter
\newlabel{algebra-lemma-completion-tensor}{{00MA}{source}{}{sourcepointer-1}{}}
\newlabel{algebra-lemma-completion-tensor@cref}{{[originallemma][0][]00MA}{[1][0][]source}{}{}{}}
\makeatother
\crefname{originalreference}{original reference}{original references}
\Crefname{originalreference}{Original reference}{Original references}
\makeatletter
\newlabel{lemma-NL-is-limit}{{0AJS}{source}{}{sourcepointer-2}{}}
\newlabel{lemma-NL-is-limit@cref}{{[originalreference][0][]0AJS}{[1][0][]source}{}{}{}}
\makeatother
\crefname{originalreference}{original reference}{original references}
\Crefname{originalreference}{Original reference}{Original references}
\makeatletter
\newlabel{lemma-exact-sequence-NL}{{00S2}{source}{}{sourcepointer-3}{}}
\newlabel{lemma-exact-sequence-NL@cref}{{[originalreference][0][]00S2}{[1][0][]source}{}{}{}}
\makeatother
\crefname{originallemma}{original lemma}{original lemmas}
\Crefname{originallemma}{Original lemma}{Original lemmas}
\makeatletter
\newlabel{more-algebra-lemma-consequence-Artin-Rees}{{0EGU}{source}{}{sourcepointer-4}{}}
\newlabel{more-algebra-lemma-consequence-Artin-Rees@cref}{{[originallemma][0][]0EGU}{[1][0][]source}{}{}{}}
\makeatother
\crefname{originallemma}{original lemma}{original lemmas}
\Crefname{originallemma}{Original lemma}{Original lemmas}
\makeatletter
\newlabel{more-algebra-lemma-transitive-lci-at-end}{{07D4}{source}{}{sourcepointer-5}{}}
\newlabel{more-algebra-lemma-transitive-lci-at-end@cref}{{[originallemma][0][]07D4}{[1][0][]source}{}{}{}}
\makeatother
\title{Fibrewise local complete intersection maps make the H\textasciicircum{}\{-1\} connecting map injective}
\author{The Stacks Project Authors}
\date{stacks-project:restricted.tex:local-snapshot-20261012; GNU FDL 1.2 or later}
\begin{document}
\maketitle
\noindent\textit{Editorial note.} All delivered mathematical text is the author's own, copied from the pinned Stacks Project chapter file restricted.tex; no statement or proof step was rewritten, reordered or removed.\par\medskip
\noindent\textit{Editorial note.} The counted proof is the author's complete argument: the proof reduces the cohomology groups to limits over the reductions via Artin-Rees completeness and lemma-NL-is-limit, applies the algebraic transitivity lemma fibrewise, and concludes injectivity of the limit map\par\medskip
\noindent\textit{Editorial note.} The label lemma-transitive-lci-at-end names the counted statement, and the proof environment immediately following it in the source is delivered verbatim as its proof.\par\medskip
\noindent\textit{Editorial note.} The counted statement is carried by the author's own lemma environment in the chapter source.\par\medskip
\noindent\textit{Allowed standard prerequisites (editorial).}\begin{itemize}
\item the exact sequence of completed naive cotangent complexes for A -> B -> C (lemma-exact-sequence-NL of the same chapter)
\item Artin-Rees consequence: H\textasciicircum{}\{-1\} of the completed complex is I-adically complete as a finite module, and equals the limit of its reductions (More on Algebra, Lemma more-algebra-lemma-consequence-Artin-Rees; Algebra, Lemma algebra-lemma-completion-tensor)
\item the identification NL\textasciicircum{}wedge\_\{B\_n/A\_n\} = NL\textasciicircum{}wedge\_\{B/A\} tensor B\_n (lemma-NL-is-limit of the same chapter)
\item the algebraic transitivity statement: fibrewise local complete intersection maps give injectivity of the connecting map on H\textasciicircum{}\{-1\} (More on Algebra, Lemma more-algebra-lemma-transitive-lci-at-end)
\end{itemize}\medskip
\ifdefined\hypertarget\hypertarget{sourcepointer-1}{}\fi
\noindent\textit{Source reference.} Original lemma 00MA: sources/stacks-chow/tags.txt, source lines 816--816. Prerequisite pointer invoked at proof line(s) 804; it names a standard fact of a sibling chapter that the counted argument uses, and every step of the counted proof built on it is retained verbatim in the delivered body.\par\medskip
\ifdefined\hypertarget\hypertarget{sourcepointer-2}{}\fi
\noindent\textit{Source reference.} Original reference 0AJS: sources/stacks-restricted/pinned-original-source.tex, source lines 596--596. Prerequisite pointer invoked at proof line(s) 806; it names a standard fact of the same chapter that the counted argument uses, and every step of the counted proof built on it is retained verbatim in the delivered body.\par\medskip
\ifdefined\hypertarget\hypertarget{sourcepointer-3}{}\fi
\noindent\textit{Source reference.} Original reference 00S2: sources/stacks-restricted/pinned-original-source.tex, source lines 725--725. Prerequisite pointer invoked at proof line(s) 781; it names a standard fact of the same chapter that the counted argument uses, and every step of the counted proof built on it is retained verbatim in the delivered body.\par\medskip
\ifdefined\hypertarget\hypertarget{sourcepointer-4}{}\fi
\noindent\textit{Source reference.} Original lemma 0EGU: sources/stacks-chow/tags.txt, source lines 18109--18109. Prerequisite pointer invoked at proof line(s) 801; it names a standard fact of a sibling chapter that the counted argument uses, and every step of the counted proof built on it is retained verbatim in the delivered body.\par\medskip
\ifdefined\hypertarget\hypertarget{sourcepointer-5}{}\fi
\noindent\textit{Source reference.} Original lemma 07D4: sources/stacks-chow/tags.txt, source lines 9399--9399. Prerequisite pointer invoked at proof line(s) 809; it names a standard fact of a sibling chapter that the counted argument uses, and every step of the counted proof built on it is retained verbatim in the delivered body.\par\medskip
\begin{lemma}
\label{lemma-transitive-lci-at-end}
With assumptions as in Lemma \ref{lemma-exact-sequence-NL}
assume that $B/I^nB \to C/I^nC$ is a local complete intersection
homomorphism for all $n$. Then
$H^{-1}(\NL_{B/A}^\wedge \otimes_B C) \to H^{-1}(\NL_{C/A}^\wedge)$
is injective.
\end{lemma}

\begin{proof}
For each $n \geq 1$ we set $A_n = A/I^n$, $B_n = B/I^nB$, and
$C_n = C/I^nC$. We have
\begin{align*}
H^{-1}(\NL_{B/A}^\wedge \otimes_B C)
& =
\lim H^{-1}(\NL_{B/A}^\wedge \otimes_B C_n) \\
& =
\lim H^{-1}(\NL_{B/A}^\wedge \otimes_B B_n \otimes_{B_n} C_n) \\
& =
\lim H^{-1}(\NL_{B_n/A_n} \otimes_{B_n} C_n)
\end{align*}
The first equality follows from
More on Algebra, Lemma \ref{more-algebra-lemma-consequence-Artin-Rees}
and the fact that $H^{-1}(\NL_{B/A}^\wedge \otimes_B C)$ is a finite
$C$-module and hence $I$-adically complete for example by
Algebra, Lemma \ref{algebra-lemma-completion-tensor}.
The second equality is trivial.
The third holds by Lemma \ref{lemma-NL-is-limit}.
The maps $H^{-1}(\NL_{B_n/A_n} \otimes_{B_n} C_n) \to
H^{-1}(\NL_{C_n/A_n})$ are injective by
More on Algebra, Lemma \ref{more-algebra-lemma-transitive-lci-at-end}.
The proof is finished because we also have
$H^{-1}(\NL_{C/A}^\wedge) = \lim H^{-1}(\NL_{C_n/A_n})$
similarly to the above.
\end{proof}

\end{document}
